\scanpage{81}
\mspage{329}

$3^\circ)$ On construit \unsure{les}
\[
h_0=\unsure{\zeta|\xi}^{-1}\ell_0^{\unsure{\gamma|\delta}}
\in\widehat{\Pi}_{0,3}
\]
La relation du $2^\circ)$ assure que $h_0\sigma_\infty(h_0)=1$ i.e.
que $h_0$ peut se mettre sous la forme
\[
h_0=\sigma_\infty(g_0)^{-1}g_0
\qquad g_0\in\widehat{\Pi}_{0,3}.
\]

\margin{NB Il faudrait s'assurer par contre que la connaissance de $g_0$
en \unsure{détermine} \illegible{1} $\unsure{\gamma|\delta}$ i.e. $p$
i.e. que \deleted{\illegible{2}} le $\unsure{\gamma|\delta}$
tel que $\unsure{\zeta|\xi}=\ell_0^{\unsure{\gamma|\delta}}h_0^{-1}$
\unsure{satisfasse} (A) \illegible{3} \unsure{est unique}.
À éclaircir!}

$4^\circ)$ On peut maintenant poser
\[
g=\sigma_\infty(g_0)\ell_1^{\unsure{\gamma|\delta}}\rho(g_0)^{-1}
=(g_0\ell_0^{-\unsure{\gamma|\delta}})\unsure{\zeta|\xi}
\rho(g_0\ell_0^{-\unsure{\gamma|\delta}})^{-1}
\]
et on \deleted{\unsure{définit}} a ce qu'il faut pour définir $u$, via
\[
\left\{
\begin{aligned}
u(\sigma_\infty)&=\sigma_\infty\\
u(\rho)&=g\rho
\end{aligned}
\right.
\]

Ici, la seule équation délicate, \unsure{\deleted{pour}} \unsure{laquelle|lesquelles}
l'existence (non l'unicité, en bon cas) d'une solution, est problématique,
est l'équation C (en $\unsure{\gamma|\delta}\in\widehat{\mathbb Z}$)
--- plus la condition encore que $2\unsure{\gamma|\delta}+1\in\widehat{\mathbb Z}^{*}$.
C'est donc une \deleted{\unsure{question}}\inserted{\uline{condition}}
sur $\lambda\in\widehat{\Pi}_{0,3}$ --- ou plutôt deux types de condition,
suivant que $\unsure{\zeta|\xi}$ est défini en termes de $\lambda$ par
l'une ou l'autre formule (B).

Je prévois qu'on est dans l'une ou

\margin{80}

\scanpage{82}

l'autre cas des formules (B), suivant que $p\bmod 3$ est congru à $1$, ou à $-1$
(NB. il ne doit pas être congru à $0$, puisque $p\bmod 3\in(\mathbb Z/3\mathbb Z)^*$)
i.e. suivant que $2\unsure{\gamma|\delta}\equiv 0$ ou $-2$ $(3)$ i.e.
$\unsure{\gamma|\delta}\equiv 0$ ou $-1$ $(3)$
(NB. il ne faut pas que $\unsure{\gamma|\delta}$ soit $\equiv 1$ $(3)$,
sinon $2\unsure{\gamma|\delta}+1=p$ serait congru à $0\bmod(3)$).
C'est en tous cas ce qui \unsure{sera} le cas si $u\in\Gamma_{\mathbb Q_0}$.

Pour le tester, réduisons le \deleted{\illegible{1}}\inserted{groupe}
$\widehat{\unsure{\mathcal G|\widetilde{\mathcal G}}}_{0,3}^{+}$
mod les \unsure{équations}
$\ell_0^3=1$ i.e. $(\sigma_\infty\rho)^6=1$
\[
\deleted{\ell_0^3=\ell_1^3=\ell_\infty^3=1\quad\illegible{1}}
\]
de sorte que $\widehat{\unsure{\mathcal G|\widetilde{\mathcal G}}}_{0,3}^{+}$
est \unsure{remplacé} par une extension de $\unsure{\mathfrak S|\mathcal S}_3$
par le groupe défini par les générateurs $\ell_0,\ell_1,\ell_\infty$ satisfaisant
\[
\unsure{\ell_0}^3=\ell_1^3=\ell_\infty^3=(\ell_\infty\ell_1\ell_0)=1
\]
--- ou mieux, utilisons
\[
\widehat{\unsure{\mathcal G|\widetilde{\mathcal G}}}_{0,3}^{+}
\simeq\unsure{\mathrm{Sl}|\mathrm{SL}}(2,\mathbb Z)^{\prime\,\wedge}
\]
et réduisons carrément dans
\[
\unsure{\mathrm{Sl}|\mathrm{SL}}(2,\mathbb Z/6\mathbb Z)'
\simeq\unsure{\mathrm{Sl}|\mathrm{SL}}(2,\mathbb Z/2\mathbb Z)
\times\unsure{\mathrm{Sl}|\mathrm{SL}}(2,\mathbb Z/3\mathbb Z)'.
\]

\margin{81}

\scanpage{83}
\mspage{330}

Mais pour faire les calculs, il faut pouvoir identifier $\sigma_\infty$
\inserted{et} $\rho$ comme éléments de
$\unsure{\mathrm{Sl}|\mathrm{SL}}(2,\unsure{\widehat{\mathbb Z}|\mathbb Z})$
--- explicitation que nous ferons tantôt.

\begin{center}\rule{2cm}{0.4pt}\end{center}

Revenant à (24 bis), et y portant
$\unsure{h_0|h_\infty}=\unsure{\zeta|\xi}^{-1}\unsure{\ell_0|\ell_\infty}^{\unsure{\gamma|\delta}}$, il vient
\[
\tag{61}
g_0=\unsure{\widehat{\varphi}|\widehat{\rho}}(\unsure{\zeta|\xi})
\varepsilon_0^{\unsure{\gamma|\delta}-\beta}
\]
On pense donc que $u$ correspond à
$(\unsure{\zeta|\xi},g_0,\unsure{\gamma|\delta},\beta)$,
avec $\unsure{\zeta|\xi},g_0\in\widehat{\Pi}_{0,3}$
et $\unsure{\gamma|\delta},\beta\in\widehat{\mathbb Z}$
\margin{\illegible{1}}
satisfaisant
\[
\tag{62}
\left\{
\begin{aligned}
&\unsure{\zeta|\xi}\rho(\unsure{\zeta|\xi})\rho^2(\unsure{\zeta|\xi})=1\\
&\unsure{\zeta|\xi}\underbrace{\sigma_\infty(g_0)^{-1}g_0}_{\unsure{h|h_0}}
=\ell_0^{\unsure{\gamma|\delta}}\\
&\deleted{\illegible{1}}g_0
=\unsure{\widehat{\varphi}|\widehat{\rho}}(\unsure{\zeta|\xi})
\varepsilon_0^{\unsure{\gamma|\delta}-\beta}\\
&2\unsure{\gamma|\delta}+1\in\widehat{\mathbb Z}^{*}
\end{aligned}
\right.
\]
dont on peut éliminer $g_0$, et on trouve des équations en
$(\unsure{\zeta|\xi},\unsure{r|\gamma},\unsure{s|\beta})$
$(r=\unsure{\gamma|\delta}-\beta,\ s=\unsure{\gamma|\delta}+\beta)$
\[
\tag{63}
\left\{
\begin{aligned}
&\unsure{\zeta|\xi}\rho(\unsure{\zeta|\xi})\rho^2(\unsure{\zeta|\xi})=1\\
&\sigma_\infty\bigl(\unsure{\widehat{\varphi}|\widehat{\rho}}(\unsure{\zeta|\xi})\bigr)^{-1}
\unsure{\widehat{\varphi}|\widehat{\rho}}(\unsure{\zeta|\xi})
=\varepsilon_1^r\unsure{\zeta|\xi}^{-1}\varepsilon_0^s\\
&\underbrace{r+s+1}_{p}\in\widehat{\mathbb Z}^{*}
\end{aligned}
\right.
\]
\margin{(\unsure{donc|avec} $\unsure{\zeta|\xi}=\lambda\rho(\lambda)^{-1}$
ou $\unsure{\zeta|\xi}=\lambda\ell_0^{-1}\rho(\lambda)^{-1}$}

Ici, $\unsure{\lambda}$ est déterminé mod transformation en
$\lambda'=\deleted{\illegible{1}}\ell_0^\alpha\lambda$
donc $\unsure{\zeta|\xi}$ devient
$\unsure{\zeta|\xi}'=\ell_0^\alpha\unsure{\zeta|\xi}\ell_1^{-\alpha}$

\margin{82}

\scanpage{84}
\mspage{331}

\uline{42} Récapitulation sur le groupe
\[
(\unsure{\mathcal G|\widetilde{\mathcal G}}_{0,3})
\simeq
\bigl(\unsure{\Pi_1|\pi_1}(U_{0,3}^{\mathrm{an}},\unsure{\mathfrak S|\mathcal S}_3;-\bar j)\bigr)
\simeq
\bigl(\widetilde\Pi'_{0,3}\simeq
\unsure{\Pi_1|\pi_1}((\mathbb P^{1\mathrm{an}},0\{0\}+2\{1\}+3\{\infty\}),\unsure{\{1,\tau\}};-\bar j)\bigr)
\]
\[
\simeq\bigl(\mathrm{Gl}(2,\mathbb Z)'=\mathrm{Gl}(2,\mathbb Z)/\{\pm1\}
\simeq\mathrm{GP}(1,\mathbb Z)\bigr)
\simeq\bigl(\mathcal T'_{1,1}=\mathcal T_{1,1}/\{\pm1\}\bigr)
\]
et sur les automorphismes ``de lacets'' de
$\widehat{\unsure{\mathcal G|\widetilde{\mathcal G}}}_{0,3}^{+}$ et de
$\widehat\Pi_{0,3}$ (Formulaire fondamental).

\margin{\deleted{\illegible{1}} $\unsure{\mathcal G|\widetilde{\mathcal G}}_{0,3}$ défini comme $\mathrm{Aut}_{\unsure{\mathrm{lac}|\mathrm{loc}}}(\Pi_{0,3})$ --- \unsure{on en donne cinq} \unsure{définitions} différentes ! (En fait, il y \illegible{3}, cf. (\unsure{V}) ci-dessous\ldots)}

\uline{I}) $\mathcal T_{0,3}$

Le groupe $\unsure{\mathcal G|\widetilde{\mathcal G}}_{0,3}$ est \unsure{présenté} comme une extension
\[
(1)\qquad
\diagram{\begin{tikzcd}[ampersand replacement=\&,column sep=small]
1 \arrow[r,dash] \& \Pi_{0,3} \arrow[r] \arrow[d,phantom,"\Vert"] \& \unsure{\mathcal G|\widetilde{\mathcal G}}_{0,3} \arrow[r] \& \mathcal T_{0,3} \arrow[r] \arrow[d,phantom,"\wr"] \& 1 \\
\& \unsure{\Pi_1|\pi_1}(U_{0,3}^{\mathrm{an}};-\bar j) \&\& \unsure{\mathfrak S|\mathcal S}_3\times\mathbb Z/2\mathbb Z \&
\end{tikzcd}}
\]
où $\unsure{\mathfrak S|\mathcal S}_3\times\mathbb Z/2\mathbb Z$ est considéré comme groupe opérant sur
$U_{0,3}^{\mathrm{an}}=\mathbb P^{1\mathrm{an}}-\{0,1,\infty\}=\mathbb C-\{0,1\}$,
\deleted{par \unsure{autom.} holomorphes}
\deleted{$\unsure{\mathfrak S|\mathcal S}_3$ opérant par permutations}
--- c'est le groupe des automorphismes conformes de $\mathbb P^{1\mathrm{an}}$ qui stabilisent $\{0,1,\infty\}$.
Le groupe \margin{\unsure{(2)}}\deleted{\unsure{$\mathfrak S_3$}}\inserted{$\unsure{\mathfrak S|\mathcal S}_3=\unsure{\mathfrak S|\mathcal S}_{\{0,1,\infty\}}$} est le sous-groupe des \unsure{automorphismes|éléments} conformes, il s'identifie au groupe des permutations de $\{0,1,\infty\}$, ses six éléments sont
\unsure{$\mathrm{id}$;|$1$;}
$\dot\sigma_0,\dot\sigma_1,\dot\sigma_\infty$ (\inserted{les trois} transpositions),
$\dot\rho,\deleted{\illegible{1}}\dot\rho^{-1}$ (permutations circulaires)
donnés par les formules explicites
\[
(2)\qquad
\left\{\begin{array}{l}
\dot\sigma_0(z)=\dfrac{z}{z-1}\inserted{\unsure{=1+\dfrac{1}{1-z}|=1+\dfrac{1}{z-1}}},\quad
\dot\sigma_1(z)=\dfrac1z,\quad
\dot\sigma_\infty(z)=1-z,\\[4pt]
\dot\rho(z)=\dfrac1{1-z},\quad
\dot\rho^{-1}(z)=1-\dfrac1z=\dfrac{z-1}{z},
\end{array}\right.
\]
avec
\[
(3)\qquad
\left\{\begin{array}{l}
\dot\sigma_0\dot\sigma_1=\dot\sigma_1\dot\sigma_\infty=\dot\sigma_\infty\dot\sigma_0=\dot\rho,\\
\dot\rho\dot\sigma_0\dot\rho^{-1}=\dot\sigma_1,\quad
\dot\rho\dot\sigma_1\dot\rho^{-1}=\dot\sigma_\infty,\quad
\dot\rho\dot\sigma_\infty\dot\rho^{-1}=\dot\sigma_0
\end{array}\right.
\]
\[
(4)\qquad
\dot\sigma_0^2=\dot\sigma_1^2=\dot\sigma_\infty^2=\dot\rho^3=1.
\]
Le groupe $\mathcal T'_{0,3}=\mathbb Z/2\mathbb Z\subset\mathcal T_{0,3}$ apparaît comme le centre,
\margin{83}

\scanpage{85}

il est engendré par la conjugaison complexe $\tau$
\[
(5)\qquad ,\quad \tau(z)=\bar z,
\]
les cinq autres éléments de $\mathcal T_{0,3}$ sont
\[
\deleted{\illegible{1}}\quad(6)\qquad
\tau\dot\sigma_0=\dot\tau_0,\quad
\tau\dot\sigma_1=\dot\tau_1,\quad
\tau\dot\sigma_\infty=\dot\tau_\infty,\quad
\tau\rho,\quad\tau\rho^{-1}=(\tau\rho)^{-1},
\]
les \unsure{antiinvolutions|anti-involutions} $\tau'_i$ $(i=0,1,\infty)$ sont particulièrement intéressantes, \unsure{données} par
\[
(7)\qquad
\unsure{\dot\tau'_0(z)}=\frac{\bar z}{\bar z-1},\quad
\unsure{\dot\tau'_1(z)}=\frac1{\bar z},\quad
\dot\tau'_\infty(z)=1-\bar z.
\]

Les points fixes \deleted{\unsure{dans $U_{0,3}$}} des \inserted{\illegible{1}} autom. conformes de \deleted{\illegible{2}} $(\mathbb P^{1\mathrm{an}},\{0,1,\infty\})$ sont les suivants
\[
(8)\qquad
\left\{\begin{array}{l}
(\mathbb P^1_{\mathbb C})^{\dot\sigma_0}=\{0,2\},\quad
(\mathbb P^1_{\mathbb C})^{\dot\sigma_1}=\{1,-1\},\quad
(\mathbb P^1_{\mathbb C})^{\dot\sigma_\infty}=\{\infty,\frac12\}\\[3pt]
\mathbb P_{\mathbb C}^{1\,\dot\rho}=\mathbb P^{1\,\dot\rho^{-1}}=\{-\bar j,-j\}
\end{array}\right.
\qquad
\begin{array}{l}
\text{où }j=\exp 2i\pi/3=\dfrac{-1+i\sqrt3}{2}\quad\text{racine primitive}\quad\text{$3^{\unsure{\mathrm{ième}|\mathrm{ème}}}$ de $1$}\\[5pt]
-\bar j=\exp 2i\pi/6=\dfrac{1+i\sqrt3}{2}\quad\text{racine primitive}\quad\text{$6^{\unsure{\mathrm{ième}|\mathrm{ème}}}$ de $1$}
\end{array}
\]

\margin{NB $\mathbb P_{\mathbb C}^{1\,\dot\rho\tau}=\mathbb P_{\mathbb C}^{1\,\dot\rho^{-1}\tau}=\varnothing$}

\[
(9)\qquad
\left\{\begin{array}{l}
\mathbb P_{\mathbb C}^{1\,\tau}=\mathbb P_{\mathbb R}^{1}\\[3pt]
\mathbb P_{\mathbb C}^{1\,\dot\tau'_0}=\{x+iy\mid(x-1)^2+y^2=1\}
\inserted{=\{z\mid z+\bar z=z\bar z\}\unsure{=\frac1z|=\frac12}}
\inserted{\unsure{=\frac12-\mathbb U}},\quad
\mathbb P_{\mathbb C}^{1\,\dot\tau'_1}=\{z\mid z\bar z=1\}=\{x+iy\mid x^2+y^2=1\}=\mathbb U\\[3pt]
\mathbb P_{\mathbb C}^{1\,\dot\tau'_\infty}=\{z\mid z+\bar z=1\}=\{x+iy\mid x=\frac12\}=\frac12+i\mathbb R;
\end{array}\right.
\]

L'image la plus suggestive est sans doute celle de la sphère de Riemann, où $\tau,\dot\tau'_0,\dot\tau'_1,\dot\tau'_\infty$ apparaissent comme les \deleted{\unsure{symétries}}\inserted{\uline{réflexions}} p.r. aux \deleted{``équateurs''}\inserted{``\unsure{gr\textsuperscript{ds}|grands} cercles''} correspondants, qui pour $\tau$ est ``l'équateur'' et pour les autres sont les grands cercles \inserted{``méridiens''}, passant par les pôles $-\bar j,-j$, qui réalisent le découpage de l'équateur en six segments circulaires, par les pts de subdivision $0,\frac12,1,2,\infty,-1$.

\figure{[Sphère en perspective : pôle supérieur étiqueté $-\bar j$, pôle inférieur $-j$ ; un équateur et plusieurs arcs méridiens joignant les pôles, avec portions pointillées. Des points sont marqués sur l'équateur, étiquetés $-1$, $0$, $\infty$, $2$, $\frac12$ et $1$. Une région à gauche est hachurée ; des flèches suivent certains arcs. Des labels intérieurs \unsure{$\dot\tau'_0$}, \unsure{$\dot\tau'_1$}, \unsure{$\dot\tau'_\infty$} et \illegible{2} se chevauchent.]}

Notons aussi les formules
\[
(10)\qquad
\left\{\begin{array}{l}
\dot\rho\dot\tau'_0\dot\rho^{-1}=\dot\tau'_1,\quad
\dot\rho\dot\tau'_1\dot\rho^{-1}=\dot\tau'_\infty,\quad
\dot\rho\dot\tau'_\infty\dot\rho^{-1}=\dot\tau'_0\\[3pt]
\dot\tau'_0\dot\tau'_1=\dot\tau'_1\dot\tau'_\infty=\dot\tau'_\infty\dot\tau'_0=\dot\rho
\end{array}\right.
\]
\margin{84}

\scanpage{86}
\mspage{332}

Les quatre grands cercles (9) découpent la sphère en 12 triangles sphériques, qui forment des domaines fondamentaux pour le groupe $\mathcal T_{0,3}\simeq\unsure{\mathfrak S|\mathcal S}_3\times\mathbb Z/2$ d'ordre 12, et sont permutés entre eux de façon simplement transitive.
\deleted{\unsure{Prenant|Prenons} un de}\inserted{(C'est \unsure{aussi|ainsi}} la subdivision barycentrique des doubles triangles sphériques \unsure{déterminés} par l'équateur $\mathbb P_{\mathbb R}^{1}$, et ses trois pts de subdivision $0,1,\infty$.)
Prenons un de ces triangles sphériques, disons $(-\bar j,0,\frac12)$, les réflexions p.r. à ses \inserted{trois} côtés, i.e.
$\tau,\dot\tau'_\infty,\dot\tau'_0$, forment un syst. de générateurs de $\mathcal T_{0,3}$, avec les relations
\[
(11)\qquad
\tau^2=\dot\tau_\infty^{\prime 2}=\dot\tau_0^{\prime 2}
=\left(\underbrace{\unsure{\tau}\dot\tau'_\infty\tau}_{\dot\sigma_\infty}\right)^2
=\left(\underbrace{\dot\tau'_0\dot\tau'_\infty}_{\dot\rho^{-1}}\right)^3
=\left(\underbrace{\tau\unsure{\dot\tau'_0}\deleted{\illegible{1}}}_{\dot\sigma_0}\right)^2=1
\]

quant à $\unsure{\mathfrak S|\mathcal S}_3=\mathcal T_{0,3}^{+}$, il est engendré par
$\tau\dot\tau'_\infty=\dot\sigma_\infty$, $\dot\tau'_0\dot\tau'_\infty=\dot\rho^{\unsure{-1|1}}$ et $\tau\dot\tau'_0=\dot\sigma_0$, avec les relations
\[
(12)\qquad
\dot\sigma_0\dot\rho\dot\sigma_\infty=\dot\sigma^2=\dot\rho^3=\dot\sigma_0^2=1.
\]

\uline{II}) Le groupe $\Pi_{0,3}$ est engendré par $\ell_0,\ell_1,\ell_\infty$, liés par la relation génératrice
\[
(13)\qquad \ell_\infty \ell_1\ell_0=1
\]
Le \deleted{groupe} sous-groupe \inserted{d'indice 2}
$\unsure{\mathfrak S|\mathcal S}'_3=(\mathcal T_{0,3})_{-\bar j}=\unsure{\mathcal T_{0,3}}$,
stabilisateur de $-\bar j$, formé des $(g,\tau^\varepsilon)\in\unsure{\mathfrak S|\mathcal S}_3\times\mathbb Z/2\mathbb Z$ tels que $\unsure{\mathrm{sg}(g)|\varepsilon(g)}=(-1)^\varepsilon$, opère sur $\Pi_{0,3}$.
Ce sous-groupe est formé de $1,\dot\tau_0,\dot\tau_1,\dot\tau_\infty,\dot\rho,\dot\rho^{-1}$, et leurs opérations sur $\Pi_{0,3}$ sont notées respectivement $1,\tau'_0,\tau'_1,\tau'_\infty,\rho,\rho^{-1}$.

\margin{\uline{NB} $\tau'_0,\tau'_1,\tau'_\infty$ \unsure{étant les|étaient} notés précédemment (dans (C)) $\unsure{\sigma_0,\sigma_1,\sigma_\infty|\bar\sigma_0,\bar\sigma_1,\bar\sigma_\infty}$.}
\margin{85}

\scanpage{87}

\deleted{\unsure{(14)}}
On identifie \deleted{\unsure{plonge}} ces éléments \deleted{au} \inserted{à des éléments du} groupe
$\mathrm{Aut}_{\mathrm{lac}}(\Pi_{0,3})\simeq\unsure{\mathcal G|\widetilde{\mathcal G}}_{0,3}$,
\deleted{\illegible{3} $\unsure{\mathcal G|\widetilde{\mathcal G}}_{0,3}$ \illegible{1}}
donc \deleted{\illegible{1}} pour $g\in\unsure{\mathfrak S|\mathcal S}'_3$, $x\in\Pi_{0,3}$, on \unsure{pose|par}
$g(x)=gxg^{-1}$; et on trouve :
\[
\tag{14}
\left\{\begin{aligned}
\rho(\ell_0)&=\rho\ell_0\rho^{-1}=\ell_1,&
\rho(\ell_1)&=\rho\ell_1\rho^{-1}=\ell_\infty,&
\rho(\ell_\infty)&=\rho\ell_\infty\rho^{-1}=\ell_0,\\
\unsure{\dot\tau'_0|\tau'_0}(\ell_0)&=\ell_0^{-1},&
\unsure{\dot\tau'_0|\tau'_0}(\ell_1)&=\ell_\infty^{-1},&
\unsure{\dot\tau'_0|\tau'_0}(\ell_\infty)&=\ell_1^{-1},\\
\tau'_1(\ell_1)&=\ell_1^{-1},&
\tau'_1(\ell_\infty)&=\ell_0^{-1},&
\tau'_1(\ell_0)&=\ell_\infty^{-1},\\
\tau'_\infty(\ell_\infty)&=\ell_\infty^{-1},&
\tau'_\infty(\ell_0)&=\ell_1^{-1},&
\tau'_\infty(\unsure{\ell_0|\ell_\infty})&=\ell_1^{-1}
\end{aligned}\right.
\]
\[
\tag{15}
:\quad\tau'_0\tau'_1=\tau'_1\tau'_\infty=\tau'_\infty\tau'_0=\rho
\qquad\text{(et bien sûr \inserted{on a} }{\tau'_0}^{2}={\tau'_1}^{2}=\unsure{{\tau'_\infty}^{2}|\tau_\infty^{2}}=\rho^3=1\text{).}
\]

\inserted{Dans} Le groupe $\unsure{\mathcal G|\widetilde{\mathcal G}}_{0,3}$ on \unsure{définit} des éléments
$\sigma_0,\sigma_1,\sigma_\infty$, au dessus de $\dot\sigma_0,\dot\sigma_1,\dot\sigma_\infty$, par
\[
\tag{16}
\left\{\begin{aligned}
\sigma_0(\ell_1)&=\ell_\infty,& \sigma_0(\ell_\infty)&=\ell_1,& \sigma_0(\ell_0)&=(\ell_1\ell_\infty)^{-1}\\
\sigma_1(\ell_\infty)&=\ell_0,& \sigma_1(\ell_0)&=\ell_\infty,& \sigma_1(\ell_1)&=(\ell_\infty\ell_0)^{-1}\\
\sigma_\infty(\ell_0)&=\unsure{\ell_1},& \sigma_\infty(\unsure{\ell_1})&=\ell_0,& \sigma_\infty(\ell_\infty)&=(\ell_0\ell_1)^{-1}
\end{aligned}\right.
\]

On a
\[
\tag{17}
\left\{\begin{aligned}
\rho\sigma_0\rho^{-1}&=\sigma_1,& \rho\sigma_1\rho^{-1}&=\sigma_\infty,& \rho\sigma_\infty\rho^{-1}&=\sigma_0\\
\rho\tau_0\rho^{-1}&=\tau_1,& \rho\tau_1\rho^{-1}&=\tau_\infty,& \rho\tau_\infty\rho^{-1}&=\tau_0
\end{aligned}\right.,
\]
\[
\tag{18}
\sigma_0\sigma_1=\rho\ell_0=\deleted{\illegible{1}}\unsure{\ell_1\rho},\quad
\sigma_1\sigma_\infty=\rho\ell_1=\ell_\infty\rho,\quad
\sigma_\infty\sigma_0=\rho\ell_\infty=\ell_0\rho^{-1}.
\]

\margin{NB $\sigma_0,\sigma_1,\sigma_\infty\in\unsure{\mathcal G|\widetilde{\mathcal G}}_{0,3}^{+}$,
$\tau'_0,\tau'_1,\tau'_\infty\in\unsure{\mathcal G|\widetilde{\mathcal G}}_{0,3}^{-}$,
$\tau_0,\tau_1,\tau_\infty\in\unsure{\mathcal G|\widetilde{\mathcal G}}_{0,3}^{-}$.}

On a ainsi des éléments de $\unsure{\mathcal G|\widetilde{\mathcal G}}_{0,3}$ au dessus \unsure{de|des}
$1,\dot\sigma_0,\dot\sigma_1,\dot\sigma_\infty,\dot\rho,\dot\rho^{-1},\dot\tau'_0,\dot\tau'_1,\dot\tau'_\infty$,
mais pas encore au dessus de $\tau$ (ce qui permet d'en avoir sur $\tau\dot\rho,\tau\dot\rho^{-1}$),
on pose
$\tau_i=\underbrace{\sigma_i}_{\text{sur }\dot\sigma_i}\underbrace{\tau'_i}_{\text{sur }\dot\tau'_i}$
\inserted{$=\tau'_i\sigma_i$}
(qui est sur $\tau$, \deleted{donc}
\deleted{\illegible{1} $\in\unsure{\mathcal G|\widetilde{\mathcal G}}_{0,3}$}
\uline{$\tau_i$ est \unsure{caractérisé}} \deleted{\illegible{1}}
par la condition d'être sur $\tau$, \inserted{\unsure{de carré $1$,}} et de commuter à $\sigma_i$ (ou encore, à $\tau'_i$)
\margin{\unsure{rappeler pourquoi on a cette} \uline{\unsure{caractérisation}}.}
\[
\tag{19}
[\tau_0,\sigma_0]=[\tau_1,\sigma_1]=[\tau_\infty,\sigma_\infty]=1
\]
\margin{$\begin{aligned}
\tau_0&=\sigma_0\tau'_0&\tau_1&=\sigma_1\tau'_1&\tau_\infty&=\sigma_\infty\tau'_\infty\\
&=\tau'_0\sigma_0&&=\tau'_1\sigma_1&&=\tau'_\infty\sigma_\infty
\end{aligned}$}
\[
\tag{20}
\tau_0^2=\tau_1^2=\tau_\infty^2=1
\qquad\text{et }\tau_0,\tau_1,\tau_\infty\text{ sur }\tau.
\]
On trouve
\[
\tag{21}
\tau_\infty\tau_1=\ell_0,\quad\tau_0\tau_\infty=\ell_1,\quad\tau_1\tau_0=\ell_\infty.
\]
Pour bien faire, il faudrait expliciter la description \unsure{topologique} (faite dans \unsure{[C]|[TC]}) de
$\sigma_0,\sigma_1,\sigma_\infty,\rho,\tau'_0,\tau'_1,\tau'_\infty,\tau_0,\tau_1,\tau_\infty$.

\margin{86}

\scanpage{88}
\mspage{333}

On a
\[
\tag{22}
\left\{\begin{aligned}
(\sigma_\infty\rho)^2&=(\rho\sigma_1)^2=\ell_0=\varepsilon_0^2,&
\varepsilon_0&\overset{\unsure{\text{dfn}|\text{déf}}}{=}\sigma_\infty\rho=\rho\sigma_1\deleted{\illegible{1}}
\quad(\mathrm{NB}\ \sigma_\infty=\rho\sigma_1\rho^{-1})\\
(\sigma_0\rho)^2&=(\rho\sigma_\infty)^2=\ell_1=\varepsilon_1^2,&
\varepsilon_1&=\unsure{\sigma_0}\rho=\rho\sigma_\infty\\
(\sigma_1\rho)^2&=(\rho\unsure{\sigma_\infty|\sigma_0})^2=\ell_\infty=\varepsilon_\infty^2,&
\varepsilon_\infty&=\sigma_1\rho=\rho\sigma_0
\end{aligned}\right.
\]

L'élément $\varepsilon_i\in\unsure{\mathcal G|\widetilde{\mathcal G}}_{0,3}$
(en fait $\varepsilon_i\in\unsure{\mathcal G|\widetilde{\mathcal G}}_{0,3}^{+}$) est au dessus de $\dot\sigma_i$,
il est \uline{caractérisé} par la condition
\[
\varepsilon_i^2=\ell_i.
\]

Dans le groupe $\unsure{\mathcal T|T}_{0,3}$ d'ordre $12$, il y a sept éléments \unsure{involutifs}
$\dot\sigma_0,\dot\sigma_1,\dot\sigma_\infty,\tau,\dot\tau'_0,\dot\tau'_1,\dot\tau'_\infty$,
\deleted{et \unsure{quatre}} \inserted{deux} éléments d'ordre $3$
$\dot\rho,\dot\rho^{-1}$, \inserted{et \uline{deux d'ordre $6$}}
$\dot\rho\tau,\dot\rho^{-1}\tau$ (en plus de l'unité).

\inserted{Sauf les deux derniers (pour lesquels ce n'est pas possible$^*$)}
Nous \unsure{avons} \unsure{relevé} chacun de ces éléments en des éléments de $\overline{\mathrm m}$ ordre de
$\unsure{\mathcal G|\widetilde{\mathcal G}}_{0,3}$.
\margin{$^*$ en effet ${\mathbb P^1}^{\dot\rho\tau}={\mathbb P^1}^{\dot\rho^{-1}\tau}=\varnothing$}
On sait que les possibilités, mod $\Pi_{0,3}$-conjugaison, \unsure{de|des} tels relèvements,
correspondent aux composantes connexes des pts fixes \unsure{dans|de} $U=U_{0,3}^{\mathrm{an}}$
des autom. \unsure{topologiques} correspondants. --- pour
$\dot\sigma_0,\dot\sigma_1,\dot\sigma_\infty,\dot\tau'_0,\dot\tau'_1,\dot\tau'_\infty$
il n'y a qu'une possibilité mod conjugaison (on a choisi
$\sigma_0,\sigma_1,\sigma_\infty,\tau_0,\tau_1,\tau_\infty$
de la façon qui a semblé le plus commode pour les calculs),
pour $\tau$ il y en a trois --- ce sont les trois choix $\tau_0,\tau_1,\tau_\infty$
--- pour $\dot\rho,\dot\rho^{-1}$ il y en a deux,
\unsure{correspondants|correspondantes} aux deux pôles $-\bar j,-j$ --- nous avons choisi celle qui
correspond au pôle $-\bar j$.
Les relèvements correspondants à $-j$ \unsure{peuvent} s'en déduire en \unsure{conjuguant} par un $\tau_i$ :

\margin{87}

\scanpage{89}

\[
(23)\quad
\begin{gathered}
\rho_0=\tau_0\rho\tau_0^{-1}
\inserted{=\ell_\infty^{-1}\rho=\unsure{\rho \ell_1^{-1}}},\qquad
\rho_1=\tau_1\rho\tau_1^{-1}
\inserted{=\ell_0^{-1}\rho=\rho \ell_\infty^{-1}},\\
\rho_\infty=\tau_\infty\rho\unsure{\tau_\infty^{-1}}
\inserted{=\ell_1^{-1}\rho=\rho \ell_0^{-1}}
\end{gathered}
\]
sont d'ordre 3, au dessus de $\dot\rho$, et $\Pi_{0,3}$-conjugués entre eux\textsuperscript{*}, mais non $\Pi_{0,3}$-conjugués à $\rho$.

\margin{* NB
\[
\begin{gathered}
\rho\rho_0\rho^{-1}=\rho_1\\
\rho\rho_1\rho^{-1}=\rho_\infty\\
\rho\rho_\infty\rho^{-1}=\rho_0
\end{gathered}
\]}

Notons aussi \deleted{que les} les formules \unsure{corollaires} de la théorie générale -- et du fait que les composantes conn. des ens. des pts fixes qui apparaissent ici sont 1-connexes --
\[
(24)\qquad
\Pi_{0,3}^{\sigma_i}=\Pi_{0,3}^{\rho}
=\Pi_{0,3}^{\tau_i}=\Pi_{0,3}^{\tau_i'}=\{1\}
\]

Si on regarde les pts critiques de $U_{0,3}^{\mathrm{an}}$, \unsure{relatifs} à l'opération de $\mathcal{T}_{0,3}$, on trouve les \deleted{\unsure{pts}\ \illegible{1}}\inserted{deux orbites} $(2,-1,\frac12)$ et $(-\bar j,-j)$, le stabilisateur d'un point de ces orbites est d'ordre 4 resp. 6.
Pour le pt 2 (opposé à 0), ce stabilisateur est le groupe $\mathrm{D}_2\simeq\mathbb{Z}/2\mathbb{Z}\times\mathbb{Z}/2\mathbb{Z}$, engendré par $\dot\sigma_0,\dot\tau_0$ (ou $\sigma_0,\tau_0'$) -- qui a été bel et bien relevé multiplicativement dans $\unsure{\mathcal{G}|\widetilde{\mathcal{G}}}_{0,3}$.
De $\unsure{\overline{\mathrm m}|=}$, pour le point $-\bar j$, c'est $\unsure{\mathfrak{S}|\mathcal{S}}_3'\inserted{\simeq\mathrm{D}_3}$, qui a été relevé également.
Par la théorie générale on sait que les classes de $\Pi_{0,3}$-conjugaison d'un tel relèvement correspondent aux pts fixes du groupe d'opérateurs \inserted{envisagé} sur $U_{0,3}$, pour $\mathrm{D}_2^{(0)}$ il n'y a qu'un seul point fixe, savoir 2, pour $\unsure{\mathfrak{S}|\mathcal{S}}_3'=\mathrm{D}_3$ il y en a deux, $-\bar j$ et $-j$, on avait choisi $-\bar j$.
On \unsure{trouve|trouvera} l'autre classe de conjugaison

\margin{88}

\scanpage{90}
\mspage{334}

en conjuguant par l'un quelconque des $\tau_i$ $(\tau_0,\tau_1,\tau_\infty)$.

\uline{III}\quad Relation avec le groupe cartographique \unsure{triangul-} non orienté,
\[
\begin{aligned}
\widetilde{\Pi}_{0,3}'
&=\Pi_1\bigl(\unsure{\mathbb{P}_{\mathbb{C}}^{1\circ}|\mathbb{P}_{\mathbb{C}}^{1\bullet}},\,
0\{0\}+2\{1\}+3\{\infty\};-\bar j\bigr)\\
&=\bigl(\Pi_{0,3}\cdot\{1,\tau\}\bigr)/(\ell_1^2,\ell_\infty^3),
\end{aligned}
\]
et \unsure{avec} $\mathrm{SL}(2,\mathbb{Z})$.

On a des \deleted{\illegible{1}} générateurs \inserted{involutifs} $\sigma_0,\sigma_1,\sigma_2$ de $\widetilde{\Pi}_{0,3}'$,
($\sigma_i$ dans un repère de carte \deleted{\unsure{fait échange}}\inserted{\unsure{change}} la \deleted{\illegible{2}}\inserted{de \unsure{dim} $i$} du repère, donc pour un biarc $\sigma_0$ et $\sigma_2$ sont resp. les \deleted{\unsure{passages}}\inserted{chgts d'orientation} \unsure{tangentielle|tangentielles} resp. transversale, $\sigma_1$ le passage au biarc ``\uline{adjacent}''
\figure{[Petit croquis : deux traits partent du voisinage d'un point inférieur gauche, l'un vers le haut avec une pointe de flèche, l'autre en oblique vers la droite avec une pointe de flèche. Plusieurs petits traits obliques et flèches se trouvent entre eux ; aucune étiquette alphabétique n'est lisible.]}
) avec les relations génératrices
\[
(25)\qquad
\left\{\begin{aligned}
\unsure{(\sigma_0')^2}&=\unsure{(\sigma_1')^2}=\unsure{(\sigma_2')^2}
=\unsure{(\sigma_2'\sigma_0')^2}=1\\
\unsure{\bigl(\underbrace{\sigma_1'\sigma_0'}_{\rho_f}\bigr)^3}&=1
\end{aligned}\right.
\]

\margin{on pourrait aussi noter $\unsure{\sigma_\infty'}$ au lieu de $\unsure{\sigma_2'}$.
On aurait mieux fait de les noter \unsure{carrément} $\unsure{r_0,r_1,r_\infty|\sigma_0,\sigma_1,\sigma_\infty}$ (au lieu de $\sigma_0,\sigma_1,\sigma_2$).}

Le sous-groupe \inserted{$\Pi_{0,3}'$} d'indice 2 \deleted{\illegible{1}}\inserted{(groupe cart. \unsure{triangul'} orienté)} est engendré par
\[
\unsure{\sigma_1'\sigma_0'}=\rho_f\inserted{\unsure{=\ell_\infty'}},\qquad
\unsure{\sigma_2'\sigma_1'}=\rho_s\inserted{\unsure{=\ell_0'}},\qquad
\unsure{\sigma_0'\sigma_2'}=\unsure{\sigma'|\sigma}\inserted{\unsure{=\ell_1'}},
\]
avec les relations génératrices
\[
(26)\qquad
\rho_f\sigma\rho_s\deleted{\illegible{1}}=1
\qquad
\sigma^2=\rho_f^3=1
\]
C'est donc le quotient de $\Pi_{0,3}$ par $\ell_1^2,\ell_\infty^3$
\[
(\unsure{27|17})\qquad
\left\{\begin{array}{rcl}
\Pi_{0,3}&\longrightarrow&\Pi_{0,3}'\\
\ell_0&\longmapsto&\ell_0'=\rho_s\\
\ell_1&\longmapsto&\ell_1'=\sigma\\
\ell_\infty&\longmapsto&\ell_\infty'=\rho_f
\end{array}\right.
\qquad,
\]
\deleted{\illegible{1}} cet homomorphisme se prolonge en

\deleted{\unsure{(28)|(27)}}

\margin{89}
